%%=============================================================================
%% Inleiding
%%====================================================\part{title}=========================

\chapter{\IfLanguageName{dutch}{Inleiding}{Introduction}}%
\label{ch:inleiding}
Elke werknemer, waar dan ook, heeft ooit eens te maken met stress. Soms kan dit zo extreem zijn dat het de gezondheid van de persoon ernstig schaadt. In deze bachelorproef wordt de focus gelegd op de IT-sector.

Op de Werelddag voor Geestelijke Gezondheid van het jaar 2025, leed er 29\% van Europese medewerkers aan stress of depressie \autocite{EUOSHA2025}. Als er enkel naar de IT'ers wordt gekeken, verdubbelt dat cijfer \autocite{ISACA2025}.

Omwille van deze redenen is het juist interessant om op zoek te gaan naar een oplossing die de werkstress (voor IT'ers) vermindert. 
%De inleiding moet de lezer net genoeg informatie verschaffen om het onderwerp te begrijpen en in te zien waarom de onderzoeksvraag de moeite waard is om te onderzoeken. In de inleiding ga je literatuurverwijzingen beperken, zodat de tekst vlot leesbaar blijft. Je kan de inleiding verder onderverdelen in secties als dit de tekst verduidelijkt. Zaken die aan bod kunnen komen in de inleiding~\autocite{Pollefliet2011}:

%\begin{itemize}
% \item context, achtergrond
%  \item afbakenen van het onderwerp
%  \item verantwoording van het onderwerp, methodologie
%  \item probleemstelling
%  \item onderzoeksdoelstelling
%  \item onderzoeksvraag
%  \item \ldots
%\end{itemize}


\section{\IfLanguageName{dutch}{Probleemstelling}{Problem Statement}}%
\label{sec:probleemstelling}

De kern van het probleem is dat de huidige taakverdeling te weinig rekening houdt met de individuele capaciteiten, werkdruk en mentale belasting van werknemers.

\vspace{1em}Om dit probleem verder af te bakenen, worden volgende deelvragen geformuleerd die betrekking hebben op het probleemdomein:

\vspace{0.5em}``Welke factoren binnen de IT-sector dragen het meest bij aan werkgerelateerde stress?''


\vspace{0.5em}``Hoe beïnvloedt een onevenwichtige taakverdeling de mentale gezondheid van IT-professionals?''

\section{\IfLanguageName{dutch}{Onderzoeksvraag}{Research question}}%
\label{sec:onderzoeksvraag}

Het probleem omtrent de huidige taakverdeling, wordt opgelost door gebruik te maken van een AI-model.
De centrale onderzoeksvraag luidt dan ook:

\vspace{0.5em}``Hoe kan artificiële intelligentie worden ingezet om de taakverdeling (binnen de IT-sector) te optimaliseren om zo enerzijds de mentale gezondheid bij werknemers te 'verbeteren' als anderzijds de werking van het bedrijf vlotter te laten verlopen?''

\vspace{1em}Natuurlijk moeten er nog andere deelvragen beantwoord worden zoals:

\vspace{0.5em}``Hoe creëer je zo een AI-model en welke resources heb je hiervoor nodig?''

\section{\IfLanguageName{dutch}{Onderzoeksdoelstelling}{Research objective}}%
\label{sec:onderzoeksdoelstelling}

Het eindresultaat is een goed onderbouwd AI-conceptmodel dat inzicht geeft in hoe een persoonlijk afgestemd takenpakket de werkdruk kan verminderen en daardoor de mentale gezondheid kan ondersteunen. 

\section{\IfLanguageName{dutch}{Opzet van deze bachelorproef}{Structure of this bachelor thesis}}%
\label{sec:opzet-bachelorproef}

% Het is gebruikelijk aan het einde van de inleiding een overzicht te
% geven van de opbouw van de rest van de tekst. Deze sectie bevat al een aanzet
% die je kan aanvullen/aanpassen in functie van je eigen tekst.

De rest van deze bachelorproef is als volgt opgebouwd:

\vspace{1em}In Hoofdstuk~\ref{ch:literatuurstudie} wordt dieper ingegaan op het concept taak-werknemer matching, werkgerelateerde stress in de IT-sector en de rol van Agile binnen taakverdeling. Daarnaast worden al bestaande AI-gestuurde systemen en technieken voor taakallocatie geanalyseerd en vergeleken.

\vspace{1em}In Hoofdstuk~\ref{ch:methodologie} wordt de gehanteerde onderzoeksaanpak toegelicht, waarbij gebruik wordt gemaakt van synthetisch gegenereerde data en een machine learning-model om een antwoord te formuleren op de onderzoeksvragen.

% TODO: Vul hier aan voor je eigen hoofstukken, één of twee zinnen per hoofdstuk

\vspace{0.5em}In Hoofdstuk~\ref{ch:requirements} worden de functionele en niet-functionele vereisten van het systeem vastgelegd.

\vspace{0.5em}In Hoofdstuk~\ref{ch:haalbaarheid} wordt de technische en organisatorische haalbaarheid van het voorgestelde systeem geëvalueerd.

\vspace{0.5em}In Hoofdstuk~\ref{ch:ontwikkeling} wordt de opbouw van het Proof of Concept besproken: van de datasetconstructie en modeltraining tot de effectieve taaktoewijzingslogica.

\vspace{0.5em}In Hoofdstuk~\ref{ch:risico} worden technische, organisatorische en ethische risico's in kaart gebracht, met bijzondere aandacht voor privacy en dataveiligheid in het kader van de AVG.

\vspace{1em}Tenslotte in Hoofdstuk~\ref{ch:conclusie} wordt het ontwikkelde systeem vergeleken met andere taakallocatiebenaderingen en worden de resultaten hiervan besproken. Ook wordt er een antwoord geformuleerd op de onderzoeksvragen.